A software quality team that attaches 90\% prediction sets to a defect classifier usually checks the promise on a random test split. In tabular data with repeated rows, this check is misleading: many test rows have an exact copy in the training data, so the check averages over easy repeats and the genuinely new cases that matter after deployment. We propose a check--fix--report routine. The check predicts from the data alone, before any model is trained, how far coverage of new cases can fall. The fix calibrates seen and new cases separately. The report states new-case coverage next to the usual figure. We tested the routine with random forests, LightGBM and linear models on \NCsNFlagTotal{} application data sets (software defects, network intrusion, phishing), \NNTotal{} benchmark data sets, controlled copy experiments, and \NScpN{} fresh data sets analysed under plans fixed in advance. On NASA mc1, random-forest sets reported \NCsMcOneReported{}\% coverage but covered \NCsMcOneNew{}\% of new modules, and \NCsMcOneFixed{}\% with the fix. On the fresh data, the check flagged all \NScpPos{} data sets whose new cases fell more than 2 percentage points short, and with the fix new-case coverage never fell more than \NWTwoWorstFixAbs{} points below the target. The check is a warning and a limit rather than a forecast: on the first fresh collection it did not rank data sets by the size of their shortfall. The routine is available as an open-source tool.
